% SPDX-License-Identifier: Apache-2.0
\section{Conclusion}
\label{sec:conclusion}

PQC migration is being managed as a schedule, and the artefacts that will decide
whether it succeeds are evidentiary. A cryptographic inventory that cannot be
audited and a formalization whose gaps are not stated are both forms of
assertion wearing the clothes of evidence.

This paper reported one project's attempt to produce evidence along two axes:
a tiered Isabelle/HOL formalization of the VWZ construction, and a
schema-validated CycloneDX 1.6 CBOM in CI. The formalization closes two tiers
with no remaining \texttt{sorry} and leaves 17 open across the repository; the
CBOM covers 9 components over 304 files and validates against the official
schema on every push.

The methodological finding is the part we would carry to the next project. Three
artefacts passed every internal check and were wrong: a build that succeeds
because the session permits it, a standard-format document that fails the
standard, and a page that never rendered. Each was caught by a reference external
to the project, and no internal check could have caught it, because each internal
check shared the assumption that was wrong.

The discipline that follows is not elaborate. Make one test fail on purpose, so
that passing means something. Validate standard formats against the standard, not
against your parser. Execute the code instead of reading it. State the open
obligations by name and the reason each is open. Prefer a missing field to an
invented one.

None of this is novel as individual practice. What we think is worth reporting is
that all three failures occurred in a project that was actively trying to be
careful, and that each was invisible until the external reference was applied.
The evidence a PQC migration needs is evidence someone else will check. Building
the external reference into the development is not a final audit step; it is the
difference between having evidence and having a claim.

\subsection{Future work}

\paragraph{Closing T2.} The remaining 7 linear-algebra obligations are
index bookkeeping, not open mathematics. The \texttt{i > 0} determinant case
requires splitting on $j = 0$ before the index shift, as recorded in the source;
the same boundary discipline likely applies to the other six.

\paragraph{Defining T3.} The 8 security and probabilistic obligations cannot be
proved until their subjects are defined. That is modelling work, and it is the
largest remaining item.

\paragraph{Key generation and the XOF.} Both are outside the current model.
Bringing them in is a prerequisite for any security statement rather than a
structural one.

\paragraph{Generalising the negative control.} A session that must fail is cheap
and catches the most dangerous class of error. It would be worth understanding
what the equivalent is for a CBOM pipeline: an artefact that a correct validator
must reject, run on every build.
